% Options for packages loaded elsewhere
\PassOptionsToPackage{unicode}{hyperref}
\PassOptionsToPackage{hyphens}{url}
\documentclass[
]{article}
\usepackage{xcolor}
\usepackage{amsmath,amssymb}
\setcounter{secnumdepth}{-\maxdimen} % remove section numbering
\usepackage{iftex}
\ifPDFTeX
  \usepackage[T1]{fontenc}
  \usepackage[utf8]{inputenc}
  \usepackage{textcomp} % provide euro and other symbols
\else % if luatex or xetex
  \usepackage{unicode-math} % this also loads fontspec
  \defaultfontfeatures{Scale=MatchLowercase}
  \defaultfontfeatures[\rmfamily]{Ligatures=TeX,Scale=1}
\fi
\usepackage{lmodern}
\ifPDFTeX\else
  % xetex/luatex font selection
\fi
% Use upquote if available, for straight quotes in verbatim environments
\IfFileExists{upquote.sty}{\usepackage{upquote}}{}
\IfFileExists{microtype.sty}{% use microtype if available
  \usepackage[]{microtype}
  \UseMicrotypeSet[protrusion]{basicmath} % disable protrusion for tt fonts
}{}
\makeatletter
\@ifundefined{KOMAClassName}{% if non-KOMA class
  \IfFileExists{parskip.sty}{%
    \usepackage{parskip}
  }{% else
    \setlength{\parindent}{0pt}
    \setlength{\parskip}{6pt plus 2pt minus 1pt}}
}{% if KOMA class
  \KOMAoptions{parskip=half}}
\makeatother
\usepackage{longtable,booktabs,array}
\newcounter{none} % for unnumbered tables
\usepackage{calc} % for calculating minipage widths
% Correct order of tables after \paragraph or \subparagraph
\usepackage{etoolbox}
\makeatletter
\patchcmd\longtable{\par}{\if@noskipsec\mbox{}\fi\par}{}{}
\makeatother
% Allow footnotes in longtable head/foot
\IfFileExists{footnotehyper.sty}{\usepackage{footnotehyper}}{\usepackage{footnote}}
\makesavenoteenv{longtable}
\ifLuaTeX
  \usepackage{luacolor}
  \usepackage[soul]{lua-ul}
\else
  \usepackage{soul}
\fi
\setlength{\emergencystretch}{3em} % prevent overfull lines
\providecommand{\tightlist}{%
  \setlength{\itemsep}{0pt}\setlength{\parskip}{0pt}}
\usepackage{bookmark}
\IfFileExists{xurl.sty}{\usepackage{xurl}}{} % add URL line breaks if available
\urlstyle{same}
\hypersetup{
  pdftitle={Analisi dei requisiti},
  hidelinks,
  pdfcreator={LaTeX via pandoc}}

\title{\protect\phantomsection\label{_gwvbjzmnw5ec}{}Analisi dei
requisiti}
\author{}
\date{}

\begin{document}
\maketitle

Slide per l'analisi:
\href{https://www.math.unipd.it/~tullio/IS-1/2025/Dispense/T05.pdf}{\ul{https://www.math.unipd.it/\textasciitilde tullio/IS-1/2025/Dispense/T05.pdf}}

Esempi:

\begin{itemize}
\item
  GammardX:
  \href{https://gammardx.github.io/Documents/Documenti\%20esterni/Analisi\%20dei\%20Requisiti\%20v2.0.0.pdf}{\ul{https://gammardx.github.io/Documents/Documenti\%20esterni/Analisi\%20dei\%20Requisiti\%20v2.0.0.pdf}}
\item
  Torchlight:
  \href{https://torchlight-swe2324.github.io/docs/pb/Documentazione\%20Esterna/analisi_requisiti_v2.0.0.pdf}{\ul{https://torchlight-swe2324.github.io/docs/pb/Documentazione\%20Esterna/analisi\_requisiti\_v2.0.0.pdf}}
\item
  Merl:
  \href{https://merlunipd.github.io/login-warrior/output_documenti/esterni/AdR_V2.pdf}{\ul{https://merlunipd.github.io/login-warrior/output\_documenti/esterni/AdR\_V2.pdf}}
\end{itemize}

{\def\LTcaptype{none} % do not increment counter
\begin{longtable}[]{@{}
  >{\raggedright\arraybackslash}p{(\linewidth - 8\tabcolsep) * \real{0.0885}}
  >{\raggedright\arraybackslash}p{(\linewidth - 8\tabcolsep) * \real{0.1336}}
  >{\raggedright\arraybackslash}p{(\linewidth - 8\tabcolsep) * \real{0.3670}}
  >{\raggedright\arraybackslash}p{(\linewidth - 8\tabcolsep) * \real{0.1964}}
  >{\raggedright\arraybackslash}p{(\linewidth - 8\tabcolsep) * \real{0.1819}}@{}}
\toprule\noalign{}
\begin{minipage}[b]{\linewidth}\raggedright
Ver.
\end{minipage} & \begin{minipage}[b]{\linewidth}\raggedright
Data
\end{minipage} & \begin{minipage}[b]{\linewidth}\raggedright
Descrizione
\end{minipage} & \begin{minipage}[b]{\linewidth}\raggedright
Autore
\end{minipage} & \begin{minipage}[b]{\linewidth}\raggedright
Verificatore
\end{minipage} \\
\begin{minipage}[b]{\linewidth}\raggedright
\end{minipage} & \begin{minipage}[b]{\linewidth}\raggedright
\end{minipage} & \begin{minipage}[b]{\linewidth}\raggedright
\end{minipage} & \begin{minipage}[b]{\linewidth}\raggedright
\end{minipage} & \begin{minipage}[b]{\linewidth}\raggedright
\end{minipage} \\
\begin{minipage}[b]{\linewidth}\raggedright
\end{minipage} & \begin{minipage}[b]{\linewidth}\raggedright
\end{minipage} & \begin{minipage}[b]{\linewidth}\raggedright
\end{minipage} & \begin{minipage}[b]{\linewidth}\raggedright
\end{minipage} & \begin{minipage}[b]{\linewidth}\raggedright
\end{minipage} \\
\begin{minipage}[b]{\linewidth}\raggedright
0.3.0
\end{minipage} & \begin{minipage}[b]{\linewidth}\raggedright
15/04/2026
\end{minipage} & \begin{minipage}[b]{\linewidth}\raggedright
Inizio scrittura della sezione ``Requisiti''
\end{minipage} & \begin{minipage}[b]{\linewidth}\raggedright
Sara Ristovic
\end{minipage} & \begin{minipage}[b]{\linewidth}\raggedright
\end{minipage} \\
\begin{minipage}[b]{\linewidth}\raggedright
0.2.0
\end{minipage} & \begin{minipage}[b]{\linewidth}\raggedright
12/04/2026
\end{minipage} & \begin{minipage}[b]{\linewidth}\raggedright
Scrittura della sezione 2
\end{minipage} & \begin{minipage}[b]{\linewidth}\raggedright
Dario Pirrone
\end{minipage} & \begin{minipage}[b]{\linewidth}\raggedright
\end{minipage} \\
\begin{minipage}[b]{\linewidth}\raggedright
0.1.0
\end{minipage} & \begin{minipage}[b]{\linewidth}\raggedright
08/04/2026
\end{minipage} & \begin{minipage}[b]{\linewidth}\raggedright
Creazione documento e scrittura della sezione 1
\end{minipage} & \begin{minipage}[b]{\linewidth}\raggedright
Dario Pirrone
\end{minipage} & \begin{minipage}[b]{\linewidth}\raggedright
\end{minipage} \\
\midrule\noalign{}
\endhead
\bottomrule\noalign{}
\endlastfoot
\end{longtable}
}

\hyperref[introduzione]{\textbf{Introduzione 3}}

\begin{quote}
\hyperref[scopo-del-documento]{Scopo del documento 3}

\hyperref[scopo-del-prodotto]{Scopo del prodotto {[}?{]} 3}

\hyperref[glossario]{Glossario 3}

\hyperref[riferimenti]{Riferimenti 3}

\hyperref[riferimenti-nominativi]{Riferimenti nominativi 3}

\hyperref[riferimenti-informativi]{Riferimenti informativi 3}
\end{quote}

\hyperref[descrizione-generale]{\textbf{Descrizione generale 4}}

\begin{quote}
\hyperref[obiettivi-del-prodotto]{Obiettivi del prodotto 4}

\hyperref[funzioni-del-prodotto]{Funzioni del prodotto 4}

\hyperref[caratteristiche-degli-utenti]{Caratteristiche degli utenti 5}

\hyperref[piattaforma-di-esecuzione]{Piattaforma di esecuzione {[}?{]}
5}
\end{quote}

\hyperref[casi-duso]{\textbf{Casi d'uso 5}}

\begin{quote}
\hyperref[introduzione-1]{Introduzione 5}

\hyperref[attori]{Attori 5}

\hyperref[elenco-dei-casi-duso]{Elenco dei casi d'uso 5}

\hyperref[agente-1]{Agente 1: 5}

\hyperref[agente-3]{Agente 3: 7}

\hyperref[agente-4]{Agente 4: 9}
\end{quote}

\hyperref[requisiti]{\textbf{Requisiti 10}}

\begin{quote}
\hyperref[requisiti-funzionali]{Requisiti funzionali 10}

\hyperref[requisiti-qualitativi]{Requisiti qualitativi 10}

\hyperref[_px8bqwk31ax4]{Requisiti di vincolo 10}

\hyperref[requisiti-di-sistema-operativo]{Requisiti di sistema operativo
{[}??{]} 10}

\hyperref[requisiti-prestazionali]{Requisiti prestazionali {[}?{]} 10}

\hyperref[requisiti-di-sicurezza]{Requisiti di sicurezza {[}?{]} 10}

\hyperref[tracciamento]{Tracciamento 10}

\hyperref[fonte---requisito]{Fonte - Requisito 10}

\hyperref[requisito---fonte]{Requisito - Fonte 10}

\hyperref[riepilogo-requisiti]{Riepilogo requisiti 10}
\end{quote}

\section{Introduzione}\label{introduzione}

\subsection{Scopo del documento}\label{scopo-del-documento}

Il presente documento, denominato \textbf{Analisi dei Requisiti}, ha
l\textquotesingle obiettivo di esporre in modo dettagliato e non ambiguo
i requisiti funzionali, qualitativi e di vincolo individuati per il
progetto \textbf{Code Guardian}. Tali requisiti sono stati delineati a
seguito di un\textquotesingle attenta analisi del capitolato e dei
successivi incontri con il proponente. Il documento si avvale
dell\textquotesingle uso di diagrammi e descrizioni dei \textbf{Casi
d\textquotesingle Uso (UC)} per illustrare le interazioni tra gli attori
esterni e il sistema, fornendo una base solida per le successive fasi di
progettazione, codifica e verifica.

\subsection{Scopo del prodotto {[}?{]}}\label{scopo-del-prodotto}

\subsection{Glossario}\label{glossario}

Per evitare ambiguità terminologiche e garantire una comprensione
univoca del documento, i termini tecnici, gli acronimi e le parole che
assumono un significato specifico nel contesto del progetto sono
riportati nel documento \emph{Glossario}. Questi termini sono
contrassegnati nel presente documento da una "G" a pedice ed è un
collegamento ipertestuale che rimanda direttamente alla definizione
corrispondente nel \emph{Glossario}. Tali parole vengono contrassegnate
solamente alla loro prima occorrenza in questo documento.

\subsection{Riferimenti}\label{riferimenti}

\subsubsection{Riferimenti nominativi}\label{riferimenti-nominativi}

\begin{itemize}
\item
  Norme di Progetto v1.0.0
\item
  Piano di Qualifica v1.0.0
\item
  \href{https://www.math.unipd.it/~tullio/IS-1/2025/Progetto/C2p.pdf}{\ul{Capitolato
  d'appalto C2 - Code Guardian}}
\item
  \href{https://www.math.unipd.it/~tullio/IS-1/2025/Dispense/PD1.pdf}{\ul{Regolamento
  del Progetto Didattico}}
\end{itemize}

\subsubsection{Riferimenti informativi}\label{riferimenti-informativi}

\begin{itemize}
\item
  Glossario v1.0.0
\item
  \href{https://www.math.unipd.it/~tullio/IS-1/2025/Dispense/T05.pdf}{\ul{Dispense
  Analisi dei Requisiti}}
\item
  \href{https://www.math.unipd.it/~rcardin/swea/2022/Diagrammi\%20Use\%20Case.pdf}{\ul{Dispense
  Casi d'Uso}}
\end{itemize}

\section{Descrizione generale}\label{descrizione-generale}

\subsection{Obiettivi del prodotto}\label{obiettivi-del-prodotto}

L\textquotesingle obiettivo principale del progetto è lo sviluppo di un
set di agenti intelligenti basati su modelli linguistici di grandi
dimensioni (\emph{LLM}), progettati per integrarsi in
un\textquotesingle architettura di orchestrazione già esistente fornita
dal proponente \emph{Var Group}. Il sistema, denominato \textbf{Code
Guardian}, mira ad automatizzare i processi di audit, manutenzione e
documentazione dei repository software, riducendo il carico di lavoro
manuale dei team di sviluppo e migliorando la qualità complessiva del
codice.

Nello specifico, il prodotto si prefigge di:

\begin{itemize}
\item
  \textbf{Garantire la coerenza documentale}: assicurando che il README
  di progetto e la documentazione tecnica (inline e API) siano
  costantemente allineati all\textquotesingle evoluzione funzionale
  della piattaforma.
\item
  \textbf{Elevare gli standard di sicurezza}: introducendo il paradigma
  di \emph{policy-as-code}, attraverso il quale gli agenti verificano
  non solo le vulnerabilità generiche OWASP, ma anche
  l\textquotesingle aderenza alle specifiche pratiche di sviluppo sicuro
  definite dal team.
\item
  \textbf{Migliorare la trasparenza verso il business}: automatizzando
  la redazione di changelog comprensibili al cliente finale, traducendo
  i progressi tecnici emersi dalle user stories in valore di business
  tangibile.
\end{itemize}

\subsection{Funzioni del prodotto}\label{funzioni-del-prodotto}

Il prodotto offre le sue funzionalità attraverso tre agenti
specializzati coordinati dall\textquotesingle orchestratore centrale:

\begin{enumerate}
\def\labelenumi{\arabic{enumi}.}
\item
  \textbf{Agente di Documentazione e README}:

  \begin{itemize}
  \item
    \textbf{Manutenzione README}: creazione e aggiornamento dinamico del
    file README di progetto in base alle modifiche apportate alla
    piattaforma.
  \item
    \textbf{Gestione documentazione tecnica}: generazione e
    aggiornamento di documentazione inline (\emph{docstring},
    \emph{JSDoc}) e delle \emph{API docs} per riflettere lo stato
    attuale del codice.
  \end{itemize}
\item
  \textbf{Agente di Security Review (Policy-as-Code)}:

  \begin{itemize}
  \item
    \textbf{Analisi vulnerabilità}: utilizzo delle API di \emph{Claude
    Code} per l\textquotesingle identificazione di vulnerabilità
    standard \emph{OWASP}.
  \item
    \textbf{Verifica policy interne}: controllo
    dell\textquotesingle adozione di pratiche di sviluppo sicuro
    specifiche del team, utilizzando i file di configurazione
    \emph{CLAUDE.md} presenti nel repository come contesto operativo.
  \end{itemize}
\item
  \textbf{Agente di Changelog da User Stories}:

  \begin{itemize}
  \item
    \textbf{Analisi Sprint}: lettura e filtraggio delle \emph{user
    stories} completate durante uno specifico sprint di sviluppo.
  \item
    \textbf{Traduzione Business-Oriented:} conversione del linguaggio
    tecnico delle task in un linguaggio chiaro e comprensibile per il
    cliente, producendo \emph{changelog} sia tecnici che di business.
  \end{itemize}
\end{enumerate}

\subsection{Caratteristiche degli
utenti}\label{caratteristiche-degli-utenti}

Il sistema è progettato per inserirsi in un flusso di lavoro di sviluppo
software moderno, interagendo con diverse figure professionali che
operano all\textquotesingle interno di un team IT o nel rapporto con il
cliente finale.

L\textquotesingle utente principale è rappresentato dallo
\textbf{sviluppatore}, il quale interagisce con il sistema
indirettamente attraverso il repository Git. Egli beneficia
dell\textquotesingle automatizzazione della documentazione tecnica e del
feedback immediato sulla sicurezza, che gli permette di correggere
eventuali vulnerabilità o difformità rispetto alle policy aziendali in
tempo reale, senza attendere revisioni manuali.

Una seconda figura chiave è il \textbf{Team Lead} o lo \textbf{Security
Auditor}. Questi utenti utilizzano il sistema come strumento di
monitoraggio e governance: definiscono le regole di "sviluppo sicuro"
all\textquotesingle interno dei file CLAUDE.md e delegano
all\textquotesingle agente di sicurezza il compito di verificarne
l\textquotesingle adozione costante su tutto il codice prodotto.

Infine, il sistema si rivolge agli \textbf{Stakeholder di business}
(come Project Manager o il cliente finale). Questi soggetti non
necessitano di competenze tecniche approfondite, ma hanno bisogno di
visibilità sull\textquotesingle andamento del progetto. Per loro,
l\textquotesingle agente di changelog funge da "traduttore",
trasformando il gergo tecnico delle user stories in report di valore che
sintetizzano i progressi funzionali raggiunti al termine di ogni sprint.

\subsection{Piattaforma di esecuzione
{[}?{]}}\label{piattaforma-di-esecuzione}

Il prodotto è concepito come un ecosistema di agenti integrati in
un\textquotesingle infrastruttura già esistente, operante principalmente
in ambiente \textbf{cloud AWS}. Essendo il sistema basato su interfacce
di programmazione e automazione, l\textquotesingle esecuzione è
garantita attraverso i principali ambienti di runtime come
\textbf{Node.js} e \textbf{Python}. L\textquotesingle interazione con il
codice sorgente avviene tramite la piattaforma \textbf{GitHub},
sfruttando le \textbf{GitHub Actions} come principale motore di
esecuzione per l\textquotesingle analisi e
l\textquotesingle aggiornamento automatico dei repository.

\section{Casi d'uso}\label{casi-duso}

\subsection{Introduzione}\label{introduzione-1}

\subsection{Attori}\label{attori}

\subsection{Elenco dei casi d'uso}\label{elenco-dei-casi-duso}

\subsubsection{Agente 1:}\label{agente-1}

UC 1: Gestione del README di Progetto

Questo caso d\textquotesingle uso riguarda il mantenimento del file
README.md principale per riflettere lo stato attuale del software.

\begin{itemize}
\item
  Attore Primario: Sviluppatore.
\item
  Attori Secondari: Orchestratore, Repository Git.
\item
  Trigger: Lo Sviluppatore effettua un commit/push di modifiche
  strutturali o richiede esplicitamente l\textquotesingle aggiornamento
  tramite l\textquotesingle Orchestratore.
\item
  Precondizioni: L\textquotesingle Agente 1 è attivo e ha i permessi di
  lettura/scrittura sul repository. Esiste un file README.md (anche
  vuoto).
\item
  Scenario Principale:

  \begin{itemize}
  \item
    L\textquotesingle Orchestratore notifica l\textquotesingle Agente 1
    della necessità di aggiornare il README.
  \item
    L\textquotesingle Agente analizza i file di configurazione (es.
    package.json, requirements.txt) e la struttura delle cartelle.
  \item
    L\textquotesingle Agente identifica le modifiche (nuove dipendenze,
    nuovi moduli, cambio versione).
  \item
    L\textquotesingle Agente rigenera le sezioni interessate del file
    README.md.
  \item
    L\textquotesingle Agente propone la modifica o sovrascrive il file
    (a seconda della configurazione).
  \end{itemize}
\item
  Post-condizioni: Il file README.md è aggiornato e coerente con
  l\textquotesingle ultima versione del codice sorgente.
\item
  Sottocasi:

  \begin{itemize}
  \item
    UC 1.1: Aggiornamento della guida all\textquotesingle installazione.
  \item
    UC 1.2: Aggiornamento della descrizione
    dell\textquotesingle architettura/moduli.
  \end{itemize}
\end{itemize}

UC 2: Generazione Documentazione Inline (Docstring/JSDoc)

L\textquotesingle agente interviene direttamente nel codice sorgente per
garantire che ogni componente sia documentato correttamente.

\begin{itemize}
\item
  Attore Primario: Sviluppatore.
\item
  Attori Secondari: Orchestratore.
\item
  Trigger: Lo Sviluppatore invia un file di codice
  all\textquotesingle Agente o l\textquotesingle Orchestratore rileva
  una funzione/classe priva di documentazione.
\item
  Precondizioni: Il file sorgente è scritto in un linguaggio supportato
  (es. JavaScript, Python).
\item
  Scenario Principale:

  \begin{itemize}
  \item
    L\textquotesingle Agente analizza il file sorgente alla ricerca di
    blocchi di codice non documentati o obsoleti.
  \item
    L\textquotesingle Agente interpreta la logica della funzione, i
    parametri di input e il valore di ritorno.
  \item
    L\textquotesingle Agente genera il blocco di commento standardizzato
    (es. JSDoc per TS/JS o Docstring per Python).
  \item
    L\textquotesingle Agente inserisce il commento sopra la firma della
    funzione/classe.
  \end{itemize}
\item
  Post-condizioni: Il codice sorgente presenta commenti inline
  aggiornati e conformi agli standard di progetto.
\item
  Sottocasi:

  \begin{itemize}
  \item
    UC 2.1: Correzione di documentazione esistente non più coerente con
    il codice (es. parametri cambiati).
  \item
    UC 2.2: Segnalazione allo Sviluppatore di logiche troppo complesse
    da documentare automaticamente.
  \end{itemize}
\end{itemize}

UC 3: Generazione Documentazione API

L\textquotesingle agente estrae le definizioni degli endpoint per creare
una guida tecnica per i consumatori delle API.

\begin{itemize}
\item
  Attore Primario: Sviluppatore.
\item
  Attori Secondari: Orchestratore.
\item
  Trigger: Lo Sviluppatore richiede la generazione della documentazione
  tecnica finale prima di un rilascio.
\item
  Precondizioni: Il sistema espone degli endpoint API definiti nel
  codice (es. tramite decorator o rotte espresse).
\item
  Scenario Principale:

  \begin{itemize}
  \item
    L\textquotesingle Agente scansiona i file dei controller e dei
    modelli dei dati.
  \item
    L\textquotesingle Agente estrae per ogni endpoint: metodo HTTP, URL,
    header, corpo della richiesta e possibili risposte.
  \item
    L\textquotesingle Agente formatta queste informazioni in un
    documento Markdown dedicato (es. API.md).
  \item
    L\textquotesingle Agente pubblica o salva il documento nella
    cartella di documentazione predefinita.
  \end{itemize}
\item
  Post-condizioni: È disponibile un documento tecnico aggiornato che
  descrive il funzionamento di tutte le API della piattaforma.
\item
  Sottocasi:

  \begin{itemize}
  \item
    UC 3.1: Validazione della conformità delle API rispetto a uno
    standard definito (es. nomi in camelCase).
  \end{itemize}
\end{itemize}

\subsubsection{\texorpdfstring{Agente 3: }{Agente 3: }}\label{agente-3}

UC 4: Esecuzione Scansione Vulnerabilità (OWASP)

L\textquotesingle agente utilizza le API di Claude Code per identificare
vulnerabilità standard nel codice sorgente.

\begin{itemize}
\item
  Attore Primario: Sviluppatore.
\item
  Attori Secondari: Orchestratore, API Claude Code.
\item
  Trigger: Lo Sviluppatore richiede una scansione di sicurezza o
  l\textquotesingle Orchestratore avvia il processo dopo un commit nel
  repository.
\item
  Precondizioni:

  \begin{itemize}
  \item
    L\textquotesingle Agente 3 è correttamente configurato con le API
    Key di Claude Code.
  \item
    Il codice sorgente è accessibile nel repository.
  \end{itemize}
\item
  Scenario Principale:

  \begin{itemize}
  \item
    L\textquotesingle Orchestratore invia i file sorgente (o il diff
    dell\textquotesingle ultimo commit) all\textquotesingle Agente 3.
  \item
    L\textquotesingle Agente 3 interroga le API di Claude Code inviando
    il codice per l\textquotesingle analisi.
  \item
    L\textquotesingle AI analizza il codice cercando pattern
    riconducibili alle vulnerabilità della Top 10 OWASP (es. SQL
    Injection, XSS, Hardcoded secrets).
  \item
    L\textquotesingle Agente 3 riceve l\textquotesingle esito
    dell\textquotesingle analisi.
  \item
    L\textquotesingle Agente 3 formatta i risultati evidenziando la
    gravità (High, Medium, Low) di ogni falla trovata.
  \end{itemize}
\item
  Post-condizioni: Viene generato un report tecnico delle vulnerabilità
  rilevate.
\item
  Sottocasi:

  \begin{itemize}
  \item
    UC 4.1: Scansione specifica su un singolo file modificato. (include)
  \item
    UC 4.2: Scansione completa dell\textquotesingle intero repository
    (Full Audit). (include)
  \end{itemize}
\end{itemize}

UC 5: Verifica Conformità "Policy-as-Code" (CLAUDE.md)

L\textquotesingle agente verifica che il codice rispetti le regole
specifiche definite dal team nel file di configurazione.

\begin{itemize}
\item
  Attore Primario: Sviluppatore.
\item
  Attori Secondari: Orchestratore, File di configurazione CLAUDE.md.
\item
  Trigger: Avvio della revisione di sicurezza (manuale o automatica).
\item
  Precondizioni: Nel repository è presente un file denominato CLAUDE.md
  contenente le linee guida di sviluppo sicuro del team.
\item
  Scenario Principale:

  \begin{itemize}
  \item
    L\textquotesingle Agente 3 legge il contenuto del file CLAUDE.md per
    estrarre il contesto e le regole (es. "usa sempre questa libreria
    per il logging", "non usare variabili globali").
  \item
    L\textquotesingle Agente analizza il codice sorgente confrontandolo
    con le direttive estratte.
  \item
    L\textquotesingle Agente identifica eventuali deviazioni dalle
    pratiche concordate dal team.
  \item
    L\textquotesingle Agente genera un elenco di violazioni della policy
    interna.
  \end{itemize}
\item
  Post-condizioni: Viene prodotto un elenco di non-conformità rispetto
  alle policy di progetto.
\item
  Sottocasi:

  \begin{itemize}
  \item
    UC 5.1: Validazione del file CLAUDE.md (verifica che il file di
    configurazione sia scritto correttamente). (extend)
  \end{itemize}
\end{itemize}

UC 6: Generazione Suggerimenti di Remediation

L\textquotesingle agente fornisce soluzioni pratiche per correggere i
problemi di sicurezza o le violazioni di policy trovate.

\begin{itemize}
\item
  Attore Primario: Sviluppatore.
\item
  Attori Secondari: API Claude Code.
\item
  Trigger: Conclusione con esito positivo (rilevamento falle) degli UC 4
  o UC 5.
\item
  Precondizioni: Sono state identificate vulnerabilità o violazioni di
  policy.
\item
  Scenario Principale:

  \begin{enumerate}
  \def\labelenumi{\arabic{enumi}.}
  \item
    Per ogni vulnerabilità o violazione trovata,
    l\textquotesingle Agente 3 richiede a Claude Code una possibile
    correzione.
  \item
    L\textquotesingle AI genera uno snippet di codice "sicuro" che
    risolve il problema.
  \item
    L\textquotesingle Agente presenta allo sviluppatore il confronto tra
    il codice insicuro e la soluzione proposta.
  \end{enumerate}
\item
  Post-condizioni: Lo sviluppatore riceve indicazioni chiare su come
  correggere il codice per renderlo conforme.
\end{itemize}

\subsubsection{\texorpdfstring{Agente 4: }{Agente 4: }}\label{agente-4}

UC 7: Recupero e Filtraggio User Stories

L\textquotesingle agente deve interfacciarsi con il sistema di gestione
dei task per raccogliere il materiale necessario al report.

\begin{itemize}
\item
  Attore Primario: Sviluppatore / Project Manager.
\item
  Attori Secondari: Orchestratore, Sistema di Task Management (es.
  GitHub Issues API, Jira API).
\item
  Trigger: Lo Sviluppatore richiede la generazione del report a
  conclusione di uno Sprint.
\item
  Precondizioni: 1. L\textquotesingle Agente ha accesso in lettura alle
  API del sistema di tracking. 2. Le User Stories sono etichettate con
  lo Sprint ID corretto e sono nello stato "Done" (Completate).
\item
  Scenario Principale:

  \begin{itemize}
  \item
    L\textquotesingle Orchestratore richiede all\textquotesingle Agente
    4 di analizzare lo Sprint corrente (es. "Sprint 5").
  \item
    L\textquotesingle Agente interroga il sistema di task management
    filtrando per ID Sprint e stato della storia.
  \item
    L\textquotesingle Agente valida i dati recuperati (verifica che ci
    siano descrizioni e titoli validi).
  \item
    L\textquotesingle Agente organizza internamente le informazioni
    (separazione tra bugfix, nuove feature e task tecnici).
  \end{itemize}
\item
  Post-condizioni: L\textquotesingle Agente dispone di un set di dati
  strutturato pronto per la trasformazione.
\item
  Sottocasi:

  \begin{itemize}
  \item
    UC 7.1: Segnalazione di User Stories prive di descrizione
    sufficiente.
  \end{itemize}
\end{itemize}

UC 8: Generazione Changelog Tecnico

L\textquotesingle agente produce un documento rivolto ai programmatori
che riassume i cambiamenti a basso livello.

\begin{itemize}
\item
  Attore Primario: Sviluppatore.
\item
  Attori Secondari: Orchestratore.
\item
  Trigger: Conclusione con successo dell\textquotesingle UC 7.
\item
  Precondizioni: Disponibilità dei dati filtrati dello Sprint.
\item
  Scenario Principale:

  \begin{enumerate}
  \def\labelenumi{\arabic{enumi}.}
  \item
    L\textquotesingle Agente aggrega le User Stories per tipologia
    tecnica (es. Backend, Frontend, DevOps).
  \item
    Per ogni storia, estrae i requisiti tecnici soddisfatti e i file
    modificati (se accessibili).
  \item
    L\textquotesingle Agente formatta il testo in linguaggio tecnico
    coerente (es. "Aggiornamento schema DB", "Refactoring API Auth").
  \item
    Produce un file (solitamente CHANGELOG.md) con link diretti ai task.
  \end{enumerate}
\item
  Post-condizioni: È disponibile un changelog dettagliato per
  l\textquotesingle uso interno del team di sviluppo.
\end{itemize}

UC 9: Generazione Changelog di Business

L\textquotesingle agente trasforma il gergo tecnico in benefici
tangibili per il cliente o l\textquotesingle utente finale.

\begin{itemize}
\item
  Attore Primario: Cliente / Stakeholder.
\item
  Attori Secondari: Orchestratore, Modello AI (LLM).
\item
  Trigger: Conclusione con successo dell\textquotesingle UC 7.
\item
  Precondizioni: Disponibilità dei dati filtrati dello Sprint.
\item
  Scenario Principale:

  \begin{itemize}
  \item
    L\textquotesingle Agente analizza le User Stories e le modifiche
    apportate.
  \item
    L\textquotesingle AI esegue una traduzione semantica: converte i
    termini tecnici in termini di valore (es. "Ottimizzazione indici DB"
    diventa "Migliorata la velocità di caricamento delle tabelle").
  \item
    L\textquotesingle Agente rimuove i dettagli tecnici irrilevanti per
    il cliente (es. bugfix minori di refactoring).
  \item
    Viene generato un report discorsivo e orientato al business (es.
    "Novità dello Sprint 5 per l\textquotesingle utente").
  \end{itemize}
\item
  Post-condizioni: Viene prodotto un documento pronto per essere
  presentato agli stakeholder non tecnici.
\item
  Sottocasi:

  \begin{itemize}
  \item
    UC 9.1: Adattamento del tono (formale/informale) in base alla
    configurazione.
  \end{itemize}
\end{itemize}

\section{Requisiti}\label{requisiti}

I requisiti riportati in questa sezione sono stati derivati dai casi
d'uso descritti nel capitolo 3 e dalle indicazioni emerse dagli incontri
con il proponente. In particolare, il perimetro del progetto è stato
chiarito come sviluppo di agenti specializzati da integrare in
un'architettura già esistente, e non come realizzazione di una nuova
architettura completa ad agenti.

\subsection{Requisiti funzionali}\label{requisiti-funzionali}

{\def\LTcaptype{none} % do not increment counter
\begin{longtable}[]{@{}
  >{\raggedright\arraybackslash}p{(\linewidth - 6\tabcolsep) * \real{0.1158}}
  >{\raggedright\arraybackslash}p{(\linewidth - 6\tabcolsep) * \real{0.1608}}
  >{\raggedright\arraybackslash}p{(\linewidth - 6\tabcolsep) * \real{0.4791}}
  >{\raggedright\arraybackslash}p{(\linewidth - 6\tabcolsep) * \real{0.2026}}@{}}
\toprule\noalign{}
\begin{minipage}[b]{\linewidth}\centering
\textbf{Codice}
\end{minipage} & \begin{minipage}[b]{\linewidth}\centering
\textbf{Priorità}
\end{minipage} & \begin{minipage}[b]{\linewidth}\centering
\textbf{Descrizione}
\end{minipage} & \begin{minipage}[b]{\linewidth}\centering
\textbf{Fonti}
\end{minipage} \\
\begin{minipage}[b]{\linewidth}\raggedright
RF.1
\end{minipage} & \begin{minipage}[b]{\linewidth}\raggedright
Obbligatorio
\end{minipage} & \begin{minipage}[b]{\linewidth}\raggedright
Il sistema deve aggiornare il file
\href{http://readme.md}{\ul{README.md}} del repository in modo coerente
con lo stato corrente del codice sorgente.
\end{minipage} & \begin{minipage}[b]{\linewidth}\raggedright
UC1
\end{minipage} \\
\begin{minipage}[b]{\linewidth}\raggedright
RF.2
\end{minipage} & \begin{minipage}[b]{\linewidth}\raggedright
Obbligatorio
\end{minipage} & \begin{minipage}[b]{\linewidth}\raggedright
Il sistema deve generare o aggiornare la documentazione inline del
codice, come docstring o JSDoc, per elementi sorgente supportati.
\end{minipage} & \begin{minipage}[b]{\linewidth}\raggedright
UC2
\end{minipage} \\
\begin{minipage}[b]{\linewidth}\raggedright
RF.3
\end{minipage} & \begin{minipage}[b]{\linewidth}\raggedright
\end{minipage} & \begin{minipage}[b]{\linewidth}\raggedright
\end{minipage} & \begin{minipage}[b]{\linewidth}\raggedright
\end{minipage} \\
\begin{minipage}[b]{\linewidth}\raggedright
RF.4
\end{minipage} & \begin{minipage}[b]{\linewidth}\raggedright
Obbligatorio
\end{minipage} & \begin{minipage}[b]{\linewidth}\raggedright
Il sistema deve eseguire un'analisi di sicurezza del codice alla ricerca
di vulnerabilità riconducibili alle categorie OWASP rilevanti.
\end{minipage} & \begin{minipage}[b]{\linewidth}\raggedright
UC4
\end{minipage} \\
\begin{minipage}[b]{\linewidth}\raggedright
RF.4.1
\end{minipage} & \begin{minipage}[b]{\linewidth}\raggedright
Obbligatorio
\end{minipage} & \begin{minipage}[b]{\linewidth}\raggedright
Il sistema deve poter eseguire una scansione di sicurezza su un singolo
file modificato.
\end{minipage} & \begin{minipage}[b]{\linewidth}\raggedright
UC4.1
\end{minipage} \\
\begin{minipage}[b]{\linewidth}\raggedright
RF.4.2
\end{minipage} & \begin{minipage}[b]{\linewidth}\raggedright
Obbligatorio
\end{minipage} & \begin{minipage}[b]{\linewidth}\raggedright
Il sistema deve poter eseguire una scansione di sicurezza sull'intero
repository.
\end{minipage} & \begin{minipage}[b]{\linewidth}\raggedright
UC4.2
\end{minipage} \\
\midrule\noalign{}
\endhead
\bottomrule\noalign{}
\endlastfoot
\end{longtable}
}

\subsection{}\label{section}

\subsection{Requisiti qualitativi}\label{requisiti-qualitativi}

{\def\LTcaptype{none} % do not increment counter
\begin{longtable}[]{@{}
  >{\raggedright\arraybackslash}p{(\linewidth - 6\tabcolsep) * \real{0.1158}}
  >{\raggedright\arraybackslash}p{(\linewidth - 6\tabcolsep) * \real{0.1608}}
  >{\raggedright\arraybackslash}p{(\linewidth - 6\tabcolsep) * \real{0.4791}}
  >{\raggedright\arraybackslash}p{(\linewidth - 6\tabcolsep) * \real{0.2026}}@{}}
\toprule\noalign{}
\begin{minipage}[b]{\linewidth}\centering
\textbf{Codice}
\end{minipage} & \begin{minipage}[b]{\linewidth}\centering
\textbf{Priorità}
\end{minipage} & \begin{minipage}[b]{\linewidth}\centering
\textbf{Descrizione}
\end{minipage} & \begin{minipage}[b]{\linewidth}\centering
\textbf{Fonti}
\end{minipage} \\
\begin{minipage}[b]{\linewidth}\raggedright
RQ.1
\end{minipage} & \begin{minipage}[b]{\linewidth}\raggedright
Obbligatorio
\end{minipage} & \begin{minipage}[b]{\linewidth}\raggedright
Gli output prodotti dagli agenti devono risultare utili al miglioramento
del codice, ad es. tramite segnalazioni, suggerimenti, frammenti di
codice o proposte di soluzione.
\end{minipage} & \begin{minipage}[b]{\linewidth}\raggedright
Verbale esterno
\end{minipage} \\
\begin{minipage}[b]{\linewidth}\raggedright
RQ.2
\end{minipage} & \begin{minipage}[b]{\linewidth}\raggedright
Obbligatorio
\end{minipage} & \begin{minipage}[b]{\linewidth}\raggedright
Devono essere definite modalità di test e validazione per gli agenti
sviluppati.
\end{minipage} & \begin{minipage}[b]{\linewidth}\raggedright
Verbale esterno
\end{minipage} \\
\midrule\noalign{}
\endhead
\bottomrule\noalign{}
\endlastfoot
\end{longtable}
}

\subsection{}\label{section-1}

\subsection{Requisiti di vincolo}\label{requisiti-di-vincolo}

{\def\LTcaptype{none} % do not increment counter
\begin{longtable}[]{@{}
  >{\raggedright\arraybackslash}p{(\linewidth - 6\tabcolsep) * \real{0.1158}}
  >{\raggedright\arraybackslash}p{(\linewidth - 6\tabcolsep) * \real{0.1608}}
  >{\raggedright\arraybackslash}p{(\linewidth - 6\tabcolsep) * \real{0.4791}}
  >{\raggedright\arraybackslash}p{(\linewidth - 6\tabcolsep) * \real{0.2026}}@{}}
\toprule\noalign{}
\begin{minipage}[b]{\linewidth}\centering
\textbf{Codice}
\end{minipage} & \begin{minipage}[b]{\linewidth}\centering
\textbf{Priorità}
\end{minipage} & \begin{minipage}[b]{\linewidth}\centering
\textbf{Descrizione}
\end{minipage} & \begin{minipage}[b]{\linewidth}\centering
\textbf{Fonti}
\end{minipage} \\
\begin{minipage}[b]{\linewidth}\raggedright
RV.1
\end{minipage} & \begin{minipage}[b]{\linewidth}\raggedright
Obbligatorio
\end{minipage} & \begin{minipage}[b]{\linewidth}\raggedright
Gli agenti sviluppati devono integrarsi in un'architettura già esistente
fornita dal proponente.
\end{minipage} & \begin{minipage}[b]{\linewidth}\raggedright
Verbale esterno
\end{minipage} \\
\begin{minipage}[b]{\linewidth}\raggedright
RV.2
\end{minipage} & \begin{minipage}[b]{\linewidth}\raggedright
Obbligatorio
\end{minipage} & \begin{minipage}[b]{\linewidth}\raggedright
Il progetto non deve prevedere l'addestramento di nuovi modelli, ma
l'integrazione di modelli o servizi tramite API.
\end{minipage} & \begin{minipage}[b]{\linewidth}\raggedright
Verbale esterno
\end{minipage} \\
\begin{minipage}[b]{\linewidth}\raggedright
RV.3
\end{minipage} & \begin{minipage}[b]{\linewidth}\raggedright
Obbligatorio
\end{minipage} & \begin{minipage}[b]{\linewidth}\raggedright
Il sistema deve poter operare sui repository di esempio e sul codice
reso disponibile dal proponente per integrazione e test.
\end{minipage} & \begin{minipage}[b]{\linewidth}\raggedright
Verbale esterno
\end{minipage} \\
\begin{minipage}[b]{\linewidth}\raggedright
RV.4
\end{minipage} & \begin{minipage}[b]{\linewidth}\raggedright
Obbligatorio
\end{minipage} & \begin{minipage}[b]{\linewidth}\raggedright
Il progetto \ul{non} prevede lo sviluppo di agenti: orchestratore, per
aggiornamento librerie, per sostituzione del codice obsoleto e per
mantenimento della copertura dei test, in quanto sono già in sviluppo
presso Var Group.
\end{minipage} & \begin{minipage}[b]{\linewidth}\raggedright
Verbale esterno
\end{minipage} \\
\midrule\noalign{}
\endhead
\bottomrule\noalign{}
\endlastfoot
\end{longtable}
}

\subsection{\texorpdfstring{Requisiti di sistema operativo {[}??{]}
}{Requisiti di sistema operativo {[}??{]} }}\label{requisiti-di-sistema-operativo}

\subsection{(facciamo un servizio web per cui non ci vogliono questi
requisiti?)}\label{facciamo-un-servizio-web-per-cui-non-ci-vogliono-questi-requisiti}

\subsection{Requisiti prestazionali
{[}?{]}}\label{requisiti-prestazionali}

\subsection{Requisiti di sicurezza}\label{requisiti-di-sicurezza}

{\def\LTcaptype{none} % do not increment counter
\begin{longtable}[]{@{}
  >{\raggedright\arraybackslash}p{(\linewidth - 6\tabcolsep) * \real{0.1158}}
  >{\raggedright\arraybackslash}p{(\linewidth - 6\tabcolsep) * \real{0.1608}}
  >{\raggedright\arraybackslash}p{(\linewidth - 6\tabcolsep) * \real{0.4791}}
  >{\raggedright\arraybackslash}p{(\linewidth - 6\tabcolsep) * \real{0.2026}}@{}}
\toprule\noalign{}
\begin{minipage}[b]{\linewidth}\centering
\textbf{Codice}
\end{minipage} & \begin{minipage}[b]{\linewidth}\centering
\textbf{Priorità}
\end{minipage} & \begin{minipage}[b]{\linewidth}\centering
\textbf{Descrizione}
\end{minipage} & \begin{minipage}[b]{\linewidth}\centering
\textbf{Fonti}
\end{minipage} \\
\begin{minipage}[b]{\linewidth}\raggedright
RS.1
\end{minipage} & \begin{minipage}[b]{\linewidth}\raggedright
Obbligatorio
\end{minipage} & \begin{minipage}[b]{\linewidth}\raggedright
Le eventuali modifiche proposte al codice o ai file del repository non
devono essere applicate senza validazione umana preventiva.
\end{minipage} & \begin{minipage}[b]{\linewidth}\raggedright
Verbale esterno
\end{minipage} \\
\begin{minipage}[b]{\linewidth}\raggedright
RS.2
\end{minipage} & \begin{minipage}[b]{\linewidth}\raggedright
Obbligatorio
\end{minipage} & \begin{minipage}[b]{\linewidth}\raggedright
I risultati dell'analisi di sicurezza devono riportare almeno il tipo di
problematica rilevata e una sua priorità o gravità.
\end{minipage} & \begin{minipage}[b]{\linewidth}\raggedright
UC4
\end{minipage} \\
\midrule\noalign{}
\endhead
\bottomrule\noalign{}
\endlastfoot
\end{longtable}
}

\subsection{Tracciamento}\label{tracciamento}

\subsubsection{Fonte - Requisito}\label{fonte---requisito}

\subsubsection{Requisito - Fonte}\label{requisito---fonte}

\subsubsection{Riepilogo requisiti}\label{riepilogo-requisiti}

\end{document}
